\section*{अध्याय \ref{ch:b3:climate-models} --- जलवायु मॉडल}

\begin{solution}{exo:b3:climate-models:1}
$A = 14.0 - 1.5 = \num{12.5}$ और $B = 0.14 - 0.05 = \num{0.09}$,
\unit{kcal.cm^{-2}} प्रति माह में। एक माह \qty{2.63e6}{s} का है, इसलिए
\qty{1}{kcal.cm^{-2}} प्रति माह $\num{4.184e7}/\num{2.63e6} =
\qty{15.91}{W/m^2}$ होता है: $A = \qty{199}{W/m^2}$, $B =
\qty{1.43}{W.m^{-2}.\celsius^{-1}}$। \qty{15}{\celsius} पर, $199 + 21.5 =
\qty{220}{W/m^2}$।
\end{solution}

\begin{solution}{exo:b3:climate-models:2}
$(238.0 - 211.9)/1.816 = \qty{14.4}{\celsius}$; 0.31 के लिए, $(234.6 -
211.9)/1.816 = \qty{12.5}{\celsius}$। एल्बिडो के हर 0.01 पर
\qty{3.4}{W/m^2} का ह्रास होता है, और $3.4/1.816 = \qty{1.9}{\celsius}$।
\end{solution}

\begin{solution}{exo:b3:climate-models:3}
$C = 1025\times3992\times60 = \qty{2.46e8}{J.m^{-2}.K^{-1}}$; इसे
एक वर्ष के \qty{3.156e7}{s} से भाग देने पर \qty{7.78}{W.yr.m^{-2}.K^{-1}}। तब $\tau
= 7.78/1.816 = \num{4.3}\text{ वर्ष}$, और $7.78/1.33 = \num{5.8}\text{ वर्ष}$।
\end{solution}

\begin{solution}{exo:b3:climate-models:4}
ECS $= 3.93/1.33 = \qty{2.95}{\celsius}$; TCR $\approx 3.93/1.95 =
\qty{2.02}{\celsius}$; अनुपात \num{0.68}।
\end{solution}

\begin{solution}{exo:b3:climate-models:5}
$Q = \qty{335}{W/m^2}$। हिमगोल: $(335\times0.38 - 211.9)/1.816 =
\qty{-46.6}{\celsius}$; उष्ण: $(335\times0.70 - 211.9)/1.816 =
\qty{12.4}{\celsius}$, जो अब भी \qty{10}{\celsius} से ऊपर है; रैंप: $335(0.54 +
0.016\,T) = 211.9 + 1.816\,T$ से $T = \qty{8.7}{\celsius}$ मिलता है। $N' =
\qty{-1.816}{W.m^{-2}.\celsius^{-1}}$ पहली दो अवस्थाओं पर, जो स्थायी हैं, और $-1.816
+ 335\times0.016 = \num{+3.54}$ रैंप पर, जो अस्थायी है।
\end{solution}

\begin{solution}{exo:b3:climate-models:6}
उष्ण शाखा वहाँ समाप्त होती है जहाँ उसका तापमान \qty{0}{\celsius} तक पहुँचता है:
$0.70\,Q = A$, $Q = \qty{302.7}{W/m^2}$, जिसका सौर स्थिरांक
\qty{1211}{W/m^2} है, जो आज से \qty{11.0}{\percent} कम है। हिमगोल पिघलता है
जब $0.38\,Q = A - 20B$, $Q = \qty{462.1}{W/m^2}$, $S_0 =
\qty{1848}{W/m^2}$। देहली $B(T^* - T_F)/(Q_0(1 - \alpha_F))$ है,
आज के तापमान $T^*$ और उस तापमान के बीच का अंतर, जिस पर
हिम फैलना शुरू होता है, उसे सूर्य के प्रकाश में बदलने पर: $1.816\times14.37/238 =
\qty{11.0}{\percent}$ यहाँ, और $1.816\times4.37/238 = \qty{3.3}{\percent}$
पाठ में।
\end{solution}

\begin{solution}{exo:b3:climate-models:7}
(क)~\qty{-45.5}{\celsius}। (ख)~$340\times0.38 + F = 211.9 - 18.2$: $F =
\qty{64.5}{W/m^2}$, दोगुने होने के प्रणोदन का 16 गुना है। (ग)~जब तक
प्रणोदन रहता है, तब ग्रह उष्ण शाखा पर कूद जाता है, $(238.0 + 64.5 -
211.9)/1.816 = \qty{50}{\celsius}$; इसे हटाने के बाद ग्रह उष्ण शाखा के साथ ठंडा होता है
उष्ण शाखा के साथ \qty{14.4}{\celsius} तक, जो आज के सूर्य पर विद्यमान है:
वह हिमगोल पर वापस नहीं लौटता।
\end{solution}

\begin{solution}{exo:b3:climate-models:8}
$\Sigma = 1.80/7.0 + 0.70/90 = \qty{0.2649}{yr^{-1}}$, $\Pi = 1.10\times0.70/
(7.0\times90) = \qty{1.222e-3}{yr^{-2}}$। मूलों से $\tau_f =
\num{3.84}\text{ वर्ष}$ और $\tau_s = \num{213}\text{ वर्ष}$; सन्निकटन
$7.0/1.80 = \num{3.89}\text{ वर्ष}$ और $90\times1.80/(1.10\times0.70) =
\num{210}\text{ वर्ष}$। $a_f = 0.597$ ठीक-ठीक, $1.10/1.80 = 0.611$
लगभग। ECS $= 3.93/1.10 = \qty{3.57}{\celsius}$; TCR $\approx
3.93/1.80 = \qty{2.18}{\celsius}$।
\end{solution}

\begin{solution}{exo:b3:climate-models:9}
हर $B + \beta = \num{5.555}$। भूमध्य रेखा: $(340\times1.221\times0.70 -
211.9 + 3.74\times12.81)/5.555 = \qty{22.8}{\celsius}$; ध्रुव:
$(340\times0.506\times0.38 - 211.9 + 47.9)/5.555 = \qty{-17.7}{\celsius}$।
परिवहन भूमध्य रेखा पर $3.74\times(22.8 - 12.8) = \qty{37}{W/m^2}$ हटाता है
और ध्रुव तक $3.74\times30.6 = \qty{114}{W/m^2}$ पहुँचाता है।
\end{solution}

\begin{solution}{exo:b3:climate-models:10}
(क)~$(1 - x^2)P_2' = 3x - 3x^3$, जिसका अवकलज $3 - 9x^2 = -6P_2$ है।
(ख)~परिवहन पद $-\mathcal D T_2(-6P_2) = 6\mathcal DT_2P_2$ है।
संतुलन $Q(1 + s_2P_2)(1 - \alpha) - A - B(T_0 + T_2P_2) - 6\mathcal
DT_2P_2 = 0$ सभी $x$ के लिए तभी सही है जब इसका स्थिर भाग और इसका $P_2$ भाग
शून्य हो, जिससे $T_0$ और $T_2$ मिलते हैं। (ग)~$T(0) - T(1) = T_2(-\tfrac12 - 1)
= -1.5\,T_2$। परिवहन के बिना $T_2 = 340\times(-0.476)\times0.70/1.816 =
\qty{-62.4}{\celsius}$: अंतर \qty{94}{\celsius} है। (घ)~$T_2 =
-55.3/1.5 = \qty{-36.9}{\celsius}$, इसलिए $B + 6\mathcal D = 113.3/36.9 =
3.07$ और $\mathcal D = \qty{0.21}{W.m^{-2}.\celsius^{-1}}$।
\end{solution}

\begin{solution}{exo:b3:climate-models:11}
(क)~$r = \tfrac12\Sigma\big(1 \pm \sqrt{1 - 4\Pi/\Sigma^2}\big)$;
$\sqrt{1 - \epsilon} \approx 1 - \epsilon/2$ लेने पर, मूल $\Sigma -
\Pi/\Sigma \approx \Sigma$ और $\Pi/\Sigma$ हैं। (ख)~$T_d = 0$ के साथ, $C\,\mathrm
dT/\mathrm dt = kt - (\gamma - \lambda)T$ है, जिसका विराम से आरंभ होने वाला हल $T =
k(t - \tau_f)/(\gamma - \lambda) + k\tau_f\mathrm e^{-t/\tau_f}/(\gamma -
\lambda)$, $\tau_f = C/(\gamma - \lambda)$। $t \gg \tau_f$ के लिए, $T \approx
kt/(\gamma - \lambda)$, और दोगुने होने पर, $kt = F_{2\times}$: TCR
$\approx F_{2\times}/(\gamma - \lambda)$। (ग)~$\Sigma = 1.95/8.1 +
0.62/110 = \qty{0.2464}{yr^{-1}}$, $\Pi = 1.33\times0.62/(8.1\times110) =
\qty{9.25e-4}{yr^{-2}}$; $\tau_f = \num{4.12}\text{ वर्ष}$ और $\tau_s =
\num{262}\text{ वर्ष}$ (लगभग $1/\Sigma = \num{4.06}\text{ वर्ष}$ और
$\Sigma/\Pi = \num{266}\text{ वर्ष}$); $4\Pi/\Sigma^2 = 0.061$।
\end{solution}

\begin{solution}{exo:b3:climate-models:12}
(क)~माध्य \qty{3.50}{\celsius}, मानक विचलन \qty{0.81}{\celsius}।
(ख)~ECS $= 1.88 + 1.86\,X$; सहसंबंध \num{0.992}; अवशिष्ट मानक
विचलन \qty{0.11}{\celsius}। (ग)~$1.88 + 1.86\times0.7 =
\qty{3.18}{\celsius}$, और उसका प्रसार $\sqrt{(1.86\times0.1)^2 + 0.11^2} =
\qty{0.22}{\celsius}$ है, जो समष्टि के प्रसार का लगभग एक-चौथाई है।
(घ)~$X$ को प्रतिपुष्टियों से जोड़ने वाली कोई भौतिक क्रियाविधि नहीं है; $X$ कई प्रेक्षणीय
राशियों को आज़माकर चुना गया है, इसलिए सहसंबंध संयोग हो सकता है; एक पूर्वाग्रह
सभी आठों मॉडलों में साझा है; और एक ही परिवार के मॉडलों को
स्वतंत्र गिना गया है।
\end{solution}

\begin{solution}{pb:b3:climate-models:1}
\textbf{1.} $Q_0 = 1360/4 = \qty{340}{W/m^2}$: ग्रह सूर्य का प्रकाश
$\pi R^2$ क्षेत्रफल की डिस्क पर ग्रहण करता है और उसे
$4\pi R^2$ क्षेत्रफल के गोले पर फैलाता है।

\textbf{2.} $340\times0.70 = \qty{238.0}{W/m^2}$; $211.9 + 1.816\times15 =
\qty{239.1}{W/m^2}$।

\textbf{3.} $(238.0 - 211.9)/1.816 = \qty{14.4}{\celsius}$।

\textbf{4.} $(129.2 - 211.9)/1.816 = \qty{-45.5}{\celsius}$, जो
\qty{-10}{\celsius} से नीचे है।

\textbf{5.} $340(0.54 + 0.016\,T) = 211.9 + 1.816\,T$: $T =
(183.6 - 211.9)/(1.816 - 5.44) = \qty{7.8}{\celsius}$।

\textbf{6.} $(238.0/\num{5.670e-8})^{1/4} = \qty{254.5}{K}$,
\qty{-18.6}{\celsius}: यह गर्म साम्य से \qty{33}{\celsius} नीचे है, जिसमें ग्रह का ग्रीनहाउस प्रभाव
$A$ और $B$ में समाहित है।

\textbf{7.} हिमगोल और गर्म साम्य पर $N' = -B = \qty{-1.816}{W.m^{-2}.\celsius^{-1}}$, दोनों स्थायी हैं;
बीच वाले साम्य पर $-1.816 + 5.44 =
\num{+3.62}$ है, जो अस्थायी है।

\textbf{8.} $\tau = 8.1/3.62 = \num{2.2}\text{ वर्ष}$।

\textbf{9.} $8.1/1.816 = \num{4.5}\text{ वर्ष}$।

\textbf{10.} $3.93/1.816 = \qty{2.2}{\celsius}$, जो आकलित सर्वोत्तम अनुमान \qty{3}{\celsius} से नीचे है और
उसकी \emph{संभावित} परास (\num{2.5} से \qty{4}{\celsius}) से भी नीचे है:
मॉडल का $B$, जो अक्षांशों के आर-पार फ़िट किया गया है,
पूरी जलवायु का प्रतिपुष्टि प्राचल नहीं है।

\textbf{11.} $\mathrm e^{-t/\tau} = 0.05$, $t = 3\tau =
\num{13}\text{ वर्ष}$ पर। वास्तविक जलवायु गहन महासागर को भी गर्म करती है,
जिसका समय-पैमाना सदियों का है।

\textbf{12.} $T_F = \qty{10}{\celsius}$ पर; $Q = (211.9 + 18.2)/0.70 =
\qty{328.7}{W/m^2}$।

\textbf{13.} $S_0 = \qty{1315}{W/m^2}$, आज से \qty{3.3}{\percent} कम है।

\textbf{14.} हिमगोल तक, $(328.7\times0.38 - 211.9)/1.816 =
\qty{-47.9}{\celsius}$।

\textbf{15.} $Q = (211.9 - 18.2)/0.38 = \qty{509.8}{W/m^2}$, $S_0 =
\qty{2039}{W/m^2}$, जो आज से \qty{50}{\percent} अधिक है।

\textbf{16.} गर्म शाखा से नीचे \qty{10}{\celsius} तक उतरते हुए,
\qty{1315}{W/m^2} पर \qty{-48}{\celsius} तक गिरावट, फिर हिमगोल शाखा पर
ऊपर चढ़ते हुए \qty{1360}{W/m^2} पर \qty{-45.5}{\celsius} तक:
ग्रह जमा रहता है। केवल \qty{2039}{W/m^2} से आगे ही यह
वापस उछलकर गर्म शाखा पर पहुँचेगा।

\textbf{17.} $\bar\alpha = 0.30\times0.9738 + 0.62\times0.0262 =
\num{0.3084}$; $\bar T = (340\times0.6916 - 211.9)/1.816 =
\qty{12.80}{\celsius}$।

\textbf{18.} $T_c = (340\times0.5643\times0.50 - 211.9 + 3.739\times12.80)/
5.555 = \qty{-12.26}{\celsius}$, जो \ang{72} पर पुनर्विश्लेषण से
\qty{3.5}{\celsius} अधिक गर्म है।

\textbf{19.} $5.555\times(-12.26) + 211.9\times3.059 = \num{580.1}$, जहाँ
$\beta/B = \num{2.059}$।

\textbf{20.} हर 1.7053, 1.7473, 1.7565, 1.7488 और
1.7362 हैं, और $Q/Q_0 = \num{1.0005}$, \num{0.9764}, \num{0.9713},
\num{0.9756} और \num{0.9827} हैं।

\textbf{21.} \ang{47.5} के पास, \num{0.971} पर। न्यूनतम के ध्रुव की ओर की हिम रेखाएँ
स्थायी हैं, $\mathrm dQ/\mathrm dx_s > 0$; भूमध्य रेखा की ओर वाली
अस्थायी हैं।

\textbf{22.} $Q/Q_0 = 1$ \ang{30} और \ang{40} के बीच भी है, \ang{30} के निकट
(अंतर्वेशन से \ang{30.2}): एक अस्थायी हिम रेखा। \ang{0} पर,
$Q/Q_0 = 580.1/(0.6106 + 0.7824)/340 = \num{1.225}$: हिमगोल आज के
सूर्य से \qty{22}{\percent} तक अधिक चमकीले किसी भी सूर्य के लिए टिकता है, इसलिए
यह आज भी संभव है।

\textbf{23.} \qty{2.9}{\percent}, जबकि शून्य-आयामी मॉडल के लिए \qty{3.3}{\percent} और
बुडिको के लिए \qty{1.6}{\percent}; दहलीज़ पर हिम लगभग \ang{48} पर ठहरती है,
बुडिको के \ang{50} के मुकाबले। अक्षांशीय मॉडल
वैश्विक मॉडल की तुलना में पहले हिमनद बनाता है, क्योंकि जहाँ हिम बनती है
वहाँ सूर्य का प्रकाश सबसे कमज़ोर होता है।

\textbf{24.} मॉडल ग्रह सौर स्थिरांक में
\textbf{\qty[text-series-to-math=true]{2.9}{\percent}} की गिरावट पर हिमनद बनता है
(गणना से \qty{2.87}{\percent})।
\end{solution}
